\section{Introduction}\label{sec:introduction}

Classical simulation of quantum circuits remains essential for algorithm development, compiler verification and
the validation of quantum hardware~\cite{preskill2018nisq,arute2019supremacy}. Exact statevector simulation stores all
$2^n$ complex amplitudes of an $n$ qubit register; in complex128 (16~bytes per amplitude) this is 4~GiB at $n=28$, 16~GiB
at $n=30$ and 64~GiB at $n=32$, so that every additional qubit doubles both memory and work. Distributing the state
over several devices is the established way to go beyond one device, both on CPU clusters and on
GPUs~\cite{smelyanskiy2016qhipster,jones2019quest,haner2017petabyte,deraedt2019massively,imamura2022mpiqulacs,bayraktar2023cuquantum}.
On FPGAs, statevector designs have mostly targeted one device~\cite{khalid2004fpga,lee2016fpga,pilch2019fpga,mahmud2019scaling,pcq2024,tiwari2026fpga30};
the Trefoil system distributes very large state vectors over several FPGAs backed by SATA disk arrays~\cite{niwase2023trefoil,wei2025qutrefoil}.

This paper studies a different point of that design space: a statevector held entirely in the high bandwidth memory
of four AMD Alveo U55C cards in one host, with complex128 arithmetic throughout. It is the multi card continuation of
the single card complex64 simulator of the companion paper~\cite{tiwari2026fpga30}, which reaches 30 qubits on one card
and contains the comparison of that design with CPU and GPU simulators. We compare two ways of dividing the work
between host and cards. In a baseline design, one kernel per card applies the local gates while every gate on a
partition qubit is computed on the host after the card's whole slice has been moved over PCIe. Our engine instead
executes every gate that does not need another card's amplitudes inside the FPGA, including diagonal and controlled
gates on partition qubits, and moves data only for the remaining gates.

The contributions, each supported by measurements on the four cards, are:
\begin{enumerate}
\item A complex128 statevector engine for the U55C with one compute unit per group of HBM pseudo channels, built
around a single $2\times2$ multiplexed primitive with a per operation slice enable, and a host runtime that distributes
the state over one, two or four cards. The engine holds \NumMaxWidthOneCard~qubits on one card and
\NumMaxWidthFourCards~qubits on four cards, verified amplitude by amplitude against an independent reference
(Section~\ref{subsec:res-correctness}).
\item Execution of diagonal, partition controlled and partition permutation gates without data movement (F1 to F3),
and a hardware specific compiler rewrite for $X$ type gates on partition qubits (F4). Ablations quantify each of
them (Section~\ref{subsec:res-ablation}).
\item A measured account of where the time goes: per phase timing, bytes moved over PCIe, scaling over one, two and
four cards, the sensitivity to the host, the board energy, and a model of the baseline design built only from
microbenchmarks that explains its timing curve (Sections~\ref{subsec:res-breakdown} to~\ref{subsec:res-model}).
\item An honest comparison with Qiskit on the same host and the same circuits (Section~\ref{subsec:res-cpu}).
\end{enumerate}
We make no claim that the FPGA system outperforms CPU or GPU statevector simulators. Section~\ref{subsec:res-cpu} shows
where Qiskit Aer on the \NumHostCores~core host is faster and where it is not; the GPU comparison of the same card type is
reported in the companion paper~\cite{tiwari2026fpga30}.
